\subsection{The actual determinant intersection}
\label{sec:geometry-height-assembly}

We retain the same arithmetic compactification and the same corrected
tangent bundles as in the point-height theorem.  Put $N=p+q$,
$f=[E:\mathbb Q]/2$, and $n_0=fpq$.  Throughout this section
$1\le q\le p$ and $3q+1\le N$.  The full flag dimension is
$d=fN(N-1)/2$; it is different from $n_0$.
The set of correction primes is the original
$S_B=S_0\cup S_{\mathrm{new}}(B)$, and
\[
 \Lambda_B=\log B+\log\operatorname{rd}(E)
              +\left|\frac{L'(1,\chi_{-d_E})}{L(1,\chi_{-d_E})}\right|.
\]
Here $-d_E$ denotes the host's original quadratic discriminant.
We take $B\ge\max(128,4fN^2)$ and $|S_0|\le B$.
First make the original cotangent modules and all their restricted
images locally free.  Then choose their determinant presentations,
coordinates and vertical factors.  No invariance of a nonfree
determinant under the preliminary modifications is used.

Write $\mathcal F_{B,\sigma}$ for these corrected tangent bundles.
The line underlying the base test is exactly
\[
       \mathcal L_B=
       \bigotimes_{\tau\in\Psi}
                    \det(\mathcal F_{B,\mathrm{sm}(\tau)})^{-1},
 \tag{HA1}
\]
where $\mathrm{sm}$ selects the original rank-$q$ half.  A constant
coefficient ideal in this identity is retained; its restriction to a
curve in a finite fibre has degree zero.  The source's coordinate
multidegree $\mu\mathrm{TwDeg}$ presents this line, including its
initial-prime and added-prime factors\cite{height-assembly-source}.

\begin{proposition}[Nefness on every finite fibre]
\label{prop:actual-finite-fibre-nefness}
The actual line $\mathcal L_B$ is nef on every finite fibre of the
chosen arithmetic model.
\end{proposition}

\begin{proof}
It suffices to prove nonnegative degree after pulling back to the
normalization of every integral projective curve in a finite fibre.
Finite extensions of its finite field do not change the sign.

Suppose first that its residue characteristic $\ell$ belongs to
$S_B$.  Put $n=\ell^{R_\ell}$, with the original $R_\ell\ge4$,
and retain the actual modules $\mathcal N_\ell$ and
$\mathcal Q_{\ell,\sigma}$ of
Section~\ref{sec:geometry-hodge-specialization}.  The chartwise
Hodge--Tate estimate proved there gives
\[
 \kappa_\ell\omega\subset\mathcal N_\ell,
 \qquad
 \kappa_\ell=\begin{cases}4,&\ell=2,\\
                         \zeta_\ell-1,&\ell>2.
             \end{cases}
\]
The algebraic integer $n/(\ell\kappa_\ell)$ is integral on the
original normal model.  For $\ell=2$ this is $2^{R_\ell-3}$;
for $\ell>2$ its valuation at every place above $\ell$ is
$R_\ell-1-1/(\ell-1)\ge0$, with the valuation normalized by
$v(\ell)=1$.  At other places $\kappa_\ell$ is a unit.
Consequently
\[
                        n\omega\subset\ell\mathcal N_\ell.
 \tag{HA2}
\]
Each original torsion differential
$\delta_t\in\mathcal N_\ell/n\omega$ therefore has a well-defined
global image in $\mathcal N_\ell/\ell\mathcal N_\ell$.
The labels here are global for a specific reason.  The original
level compactification has function field $\Omega$ and its
\texttt{FullTorsion} clause makes every geometric $n$-torsion
point $\Omega$-rational.  The generic-point clause identifies
$\Omega$ with the generic stalk; it is not merely a field of
constants containing roots of unity.  The construction with
$S_0\cup S_{\mathrm{new}}(B)$ retains this clause for every prime
in that union.  The local extensions of each original Weil
character, proved in the Hodge-specialization theorem, agree on
overlaps by uniqueness.  Their differentials consequently give
the asserted global sections\cite{height-assembly-source}.
These finitely many images generate it: locally $\mathcal N_\ell$
is generated by lifts of the $\delta_t$ and $n\omega$, and the
second summand disappears by (HA2).  Since
$\mathcal N_\ell\twoheadrightarrow\mathcal Q_{\ell,\sigma}$
is the original image map, its reduction is a surjection too.
Thus $\mathcal Q_{\ell,\sigma}/\ell\mathcal Q_{\ell,\sigma}$
is globally generated.  It is locally free by the chosen model
condition.  Its determinant has nonnegative degree on every such
curve.  The inverse determinant in (HA1) is precisely this determinant,
up to the prescribed constant different ideal.  Summing over the
selected labels proves the assertion at $\ell\in S_B$.

Now suppose $\ell\notin S_B$.  In particular
$\ell\nmid\operatorname{disc}(E)$ and the correction is absent.
The coefficient idempotents are integral after the existing splitting
base, so $\operatorname{Lie}(G)=\bigoplus_\sigma\mathcal T_\sigma$
and $\mathcal F_{B,\sigma}=\mathcal T_\sigma$ near this fibre.
The integral Hodge and conjugate filtrations used in
Theorem~\ref{thm:actual-flat-degree} give
\[
       \deg_C\det\mathcal T_\tau
           =\deg_C\det\mathcal T_{\bar\tau}
                         \quad(\tau\in\Psi).
 \tag{HA3}
\]
This is the finite-fibre character-determinant conclusion of that
proof, not an inference from the complex flat metric.  More explicitly,
the actual first-order extension has character determinant
\[
 \mathcal D_\tau=\det\mathcal T_\tau\otimes
                       (\det\mathcal T_{\bar\tau})^{-1}.
\]
Its Hodge and conjugate filtrations are those of
Theorem~\ref{geo-biextension-conjugate-filtration}.
Proposition~\ref{geo-frobenius-determinant-torsion} iterates their
Frobenius determinant identity around the finite coefficient orbit:
over a splitting field $\mathbb F_{\ell^a}$ it gives
$\mathcal D_\tau^{\otimes(\ell^a-1)}\simeq\mathcal O$.
It is this ratio, not either Hodge determinant separately, that has
degree zero\cite{height-assembly-determinant-record}.

The full Hodge determinant $\lambda_G=\det\omega_G$ is nef on
this fibre.  Here its extension is important.  The actual local
model is the normalization model with its original singleton
symplectic lattice, mapping to an auxiliary Siegel toroidal model;
the universal semiabelian family pulls back to the original family.
The Siegel Hodge line has a positive tensor power equal to the
pullback of an ample line on the minimal compactification
\cite[Theorem 7.2.4.1(2)--(3)]{height-assembly-lan}.
The normalization construction and its morphism to that auxiliary
model are \cite[(7.7)--(7.8), Proposition 7.11]{height-assembly-lan-ramified}.
The global dominating model and the subsequent vertical modifications
pull back the same family, so the equality of Hodge lines persists.
It holds along the boundary as well.  This proves the claimed
nefness of $\lambda_G$ without an assertion about every Hodge
subbundle in positive characteristic.

Finally (HA3) and the eigenspace decomposition give the exact identity
\[
             2\deg_C\mathcal L_B=\deg_C\lambda_G\ge0.
 \tag{HA4}
\]
This proves the proposition at every prime.  The general
characteristic-dependent nefness theorem for all flag weights is
not required for this particular determinant.
\end{proof}

\begin{proposition}[The original cusp potentials and their smoothing]
\label{prop:actual-cusp-smoothing}
Let $\varphi$ be the original logarithmic determinant metric relative
to a smooth reference metric, and let
$\psi=\varphi+U/2$ on the interior.  On the original smooth compact
complex model, $\varphi$ is smooth off the normal-crossing boundary
$D$, bounded above, and satisfies
\[
       \varphi\ge-C\log r-C,
       \qquad r=O(1)+\sum_{a}\log|z_a|^{-2}\ge1.
 \tag{HA5}
\]
The original $\psi$ extends to a bounded upper-semicontinuous
potential.  Both original metric potentials are plurisubharmonic
in holomorphic local frames.  There exist the original circle-invariant
chart smoothings and regularized maxima, with actual overlap errors
tending to zero, such that
\[
 \varphi_t\ge\varphi,\quad \psi_t\ge\psi,\quad
 \theta+dd^c\varphi_t\ge-\epsilon_t\omega_A,\quad
 \theta+dd^c\psi_t\ge-\epsilon'_t\omega_A,
 \qquad\epsilon_t,\epsilon'_t\longrightarrow0.
 \tag{HA6}
\]
Here $\theta$ is the reference curvature and $\omega_A$ is a fixed
positive reference form.  The first family converges locally in
$C^2$ off $D$, and the second pointwise and in $L^1$.  The bounded
family is invariant under complex conjugation, including its
original coordinate and vertical metric factors.
\end{proposition}

\begin{proof}
\emph{The actual local potentials.}
The period charts and the determinant identity of the point-height
theorem give, in a nonvanishing local frame,
\[
 \log h=\operatorname{Re}G+\tfrac12\log\det H(b,b),
 \qquad
 U=\max\bigl(0,\sup_I-\log\mathrm{nz}(I)\bigr).
 \tag{HA7}
\]
The determinant denotes the product over the original selected CM
labels; $G$ is holomorphic.  The supremum runs over the original
nonzero saturated isotropic coefficient sublattices.  Its index
normalization is the same as in the point-height theorem.

Here are direct reasons for the positivity and boundary bounds
needed below.  In a simply connected interior chart the polarized
period matrix is holomorphic, with imaginary part $Y>0$, and its
tangent metric is $Y^{-1}$ after a holomorphic change of frame.
For a holomorphic linearly independent frame $b$, the function
$\log\det(b^*Y^{-1}b)$ is plurisubharmonic.  This follows by
differentiating $Y^{-1}$, using $\partial\bar\partial Y=0$.
For an explicit check in one complex direction, make constant
changes of frame so $Y=1$ and $b=\binom{1}{0}$ at the point;
a holomorphic change of basis of $b$ makes
$\partial b=\binom{0}{C}$ there.  Write
$A=\partial Y=(A_{ij})$ in the same two blocks.  The Hessian
of this logarithm is
\[
       \|C-A_{21}\|_{\mathrm{HS}}^2
          +\|A_{11}\|_{\mathrm{HS}}^2
          +\|A_{12}\|_{\mathrm{HS}}^2\ge0.
\]
Since the direction was arbitrary this proves the assertion.

For a fixed rational isotropic subspace, complete its real period
basis to a symplectic basis.  Taking the quotient metric by its
period vectors is the Schur complement of their Gram block.  In
this basis it is the inverse of the imaginary part of the remaining
principal period block.  The same calculation therefore proves that
\[
          \log\det H(b,b)-\log\mathrm{nz}(I)
 \tag{HA8}
\]
is plurisubharmonic.  All index and trace factors in (HA8) are
constant in the period parameter.  This argument also applies
after restriction to the coefficient eigenspaces, which are
orthogonal and holomorphic; the isotropic coefficient subspace
and its symplectic quotient retain that action.

Only finitely many $I$ can have $\mathrm{nz}(I)\le1$ in a
fixed relatively compact interior chart.  Indeed the metrics
there are uniformly comparable.  The trace determinant formula
identifies $\mathrm{nz}(I)$, up to a fixed positive factor and
fixed positive power for each possible rank, with the norm of
the primitive exterior vector of its saturated integral sublattice.
In each exterior power of the fixed lattice,
a bounded set contains finitely many primitive integral vectors;
each such vector determines its rational subspace and saturation.
The formula for $U$ is consequently a finite maximum locally.
Equations (HA7)--(HA8) prove positivity and continuity off $D$.

Near the boundary use the exact original period form
\[
 \tau=T_0+\sum_a\frac{\log z_a}{2\pi i}B_a,
                    \qquad B_a\ge0,
 \tag{HA9}
\]
from the point-height proof.  On a sector with bounded arguments,
$\operatorname{Re}\tau$ is bounded and
$c\mathbf1\le\operatorname{Im}\tau\le Cr\mathbf1$ after
shrinking the polydisc.  To see the first inequality when
$\operatorname{Im}T_0$ is not positive on the whole space, split
off $\bigcap_a\ker B_a$: it is positive there, while the sum of
the other positive blocks tends to infinity on its complement.
A block completion of squares absorbs the bounded cross terms.
Thus, for real period coordinates $(u,v)$,
\[
 (u+\operatorname{Re}\tau v)^TY^{-1}
       (u+\operatorname{Re}\tau v)+v^TYv
                  \ge\frac c r(|u|^2+|v|^2).
 \tag{HA10}
\]
There is a matching upper bound $Cr(|u|^2+|v|^2)$.
The holomorphic changes of frame and their inverses stay bounded.
Consequently $|\varphi|\le C\log r+C$ and
$0\le U\le C\log r+C$: for the latter use (HA10) in each
exterior power and the positive minimum norm of a nonzero integral
Pl\"ucker vector.  These constants are allowed to depend on this
fixed datum; none enters the final intersection bound.

A plurisubharmonic function with upper growth $C\log r+C$ on a
punctured polydisc is bounded above on a smaller polydisc.  For
completeness, add $\eta\sum_a\log|z_a|$, $\eta>0$, and apply
the maximum principle successively in the punctured coordinates.
At an inner boundary the added negative linear logarithm dominates
$C\log r$; the outer distinguished boundary is a compact subset
of the interior.  Let the inner radii tend to zero and then
$\eta\downarrow0$.  This proves the assertion uniformly on a
smaller tangential compact set.  Apply it to both potentials in
(HA7).  The standard removable-singularity extension is their
upper-semicontinuous limit from the interior.

The corrected potential is also bounded below.  At a boundary
point take its actual saturated toric period sublattice $I$.
The exact toric Gram cancellation proved in the point-height
theorem makes (HA8) converge to a finite real value, along all
approaches to that point.  Since $U\ge-\log\mathrm{nz}(I)$,
it gives a lower bound for $\psi$.  This continuity assertion is
the analytic period and quotient-metric calculation from that
proof; it holds at every complex boundary point, independently
of whether the point is algebraic.  Compactness now makes $\psi$
bounded.  It gives (HA5) for $\varphi$ as well.

\emph{Uniform comparison of the actual local convolutions.}
Take finitely many holomorphic polydisc charts, with nested inner
and middle sets whose closures lie in the charts, and smooth
local reference potentials $q_j$.  On a common fixed neighbourhood
of every middle set, chart transitions and gauge differences have
bounded first derivatives.  Use the original nonnegative smooth
kernel $\rho$, of integral one and compact support in the unit
polydisc, invariant under complex rotations and conjugation.
Such a kernel is obtained by averaging a nonnegative bump over
the compact rotations and conjugation and normalizing its integral.

We give the comparison because boundedness alone does not assert
uniform approximation to an upper-semicontinuous function.  For
bounded plurisubharmonic $u$, let
$M_x(s)=\sup_{|z-x|_\infty\le s}u(z)$.  The maximum principle
on complex annuli shows that $M_x(e^s)$ is convex in $s$.
If $|u|\le B_0$ on the fixed polydisc of radius $\delta$, then,
for $L\ge2$ and sufficiently small $t$,
\[
 M_x(Lt)-M_x\bigl(t/(4L)\bigr)
 \le\frac{2B_0\log(4L^2)}{\log(\delta/(Lt))}
 \le\frac{4B_0\log(4L^2)}{\log(1/t)}.
 \tag{HA11}
\]
The first inequality is the adjacent-secant-slope inequality for
that convex function.

If $b_\rho=\sup\rho$ and
$c_\rho=b_\rho\pi^{n_0}(3/2)^{2n_0}$, the submean inequality
at a maximizer on the polydisc of radius $t/4$ gives
\[
  M_x(2t)-M_tu(x)
       \le c_\rho\bigl(M_x(2t)-M_x(t/4)\bigr).
 \tag{HA12}
\]
Indeed the polydisc of radius $3t/2$ about that maximizer contains
the kernel support and lies in the polydisc of radius $2t$ about
$x$.  The integral of the nonnegative function $M_x(2t)-u$ over
it is at most its volume times $M_x(2t)-M_x(t/4)$.  Multiplication
by the bound $b_\rho t^{-2n_0}$ for the kernel proves (HA12).
The reverse bound $M_tu(x)\le M_x(t)$ is immediate.

Apply (HA11)--(HA12) to $u_j=q_j+\psi$.  A change from chart
$j$ to $k$ sends a polydisc of radius $t$ into one of radius
$Lt$; its smooth gauge difference changes by at most $Lt$.
Subtracting the value of that gauge at the centre therefore gives
\[
 |M_tu_j-q_j-(M_tu_k-q_k)|
                 \le\frac C{\log(1/t)}+Ct
 \tag{HA13}
\]
on every middle-set overlap, with one $C$ for the finite atlas.
This is the required original overlap estimate.  It does not
assert that $M_tu_j-q_j$ converges uniformly to $\psi$.

For $\varphi$ use the original local truncations
$u_{j,K}=\max(q_j+\varphi,-K)$ with their removable
plurisubharmonic extensions.  Choose
$K(t)=A\log\log(1/t)$ with $A$ sufficiently large.
Every polydisc of radius $ct$ contains a point whose normal
coordinates have modulus at least $c't$.  Equation (HA5) thus gives
$\sup(q_j+\varphi)\ge-C_0\log\log(1/t)-C_1$ there.
Choosing $A>C_0$ absorbs the bounded gauge differences.  In the
transfer of maxima between charts, the added constant branch
$-K(t)$ is then below the transferred maximum.  The preceding
comparison applies without a constant gauge error from truncation.
Its oscillation bound is $O(K(t)+1)$, so the exact resulting error is
\[
 |M_tu_{j,K(t)}-q_j-(M_tu_{k,K(t)}-q_k)|
         \le C\frac{\log\log(1/t)}{\log(1/t)}+Ct.
 \tag{HA14}
\]
These are precisely the two original chart constructions and
their admissible overlap rates\cite{height-assembly-regularization-source}.

Choose each error at least $t$, and use the original regularized
maximum of the local expressions plus $4\epsilon(t)\eta_j$,
where $-1\le\eta_j\le0$, $\eta_j=0$ on the inner set, and
$\eta_j=-1$ near the edge of its middle set.  The overlap estimates
make the edge terms inactive.  The output is smooth.  The convex,
nondecreasing regularized maximum has derivative sum one, so its
curvature is bounded below by $-C\epsilon(t)\omega_A$.
An inner chart and the submean inequality prove domination of the
original potential.  At every point the bounded construction is
between its inner-chart convolution and that convolution plus
$2\epsilon(t)$; this proves its pointwise and $L^1$ convergence.
For the unbounded construction the truncation is inactive on each
compact subset off $D$ once $t$ is small.  There convolution
converges smoothly.  Differentiating the regularized maximum
twice adds errors $O(\epsilon(t)+t^2+t^4/\epsilon(t))$, which
tend to zero because $\epsilon(t)\ge t$.  This gives the asserted
local $C^2$ convergence and the uniform upper bound.

Finally include every chart with its conjugate chart, choose the
reference metric invariant under conjugation, and pair the cutoffs.
The original cusp norm and determinant metric are conjugation
invariant.  Change of variables in the kernel integral and
permutation invariance of the regularized maximum give exact
invariance of $\psi_t$.  Each original vertical factor has this
invariance too, as the norm of its rational arithmetic presentation.
The chosen smooth representative consequently agrees with the
original metric on the entire original complex model; no metric
averaging changes its values.
\end{proof}

\begin{theorem}[The original signed base intersection]
\label{thm:actual-signed-base-intersection}
For every original component $Z$, let $A_{B,Z}$ be the limiting
arithmetic top intersection of $\mathcal L_B$ with its original
determinant metric, and let $d_{B,Z}$ be its geometric top degree.
Use the original height and field normalization in both quantities.
Then these limits exist, $d_{B,Z}>0$, and
\[
 \frac{A_{B,Z}}{(n_0+1)d_{B,Z}}
                \ge-\frac{107}{2}fq\Lambda_B.
 \tag{HA15}
\]
For the actual full-flag test polynomial
\[
 \mu=\sum_{\tau}\sum_{i\ge p}X_{\tau i},\qquad
 V=\prod_\tau\prod_{\substack{a<b\\a,b<p\ \mathrm{or}\ a,b\ge p}}
                              (X_{\tau b}-X_{\tau a}),
 \qquad F_0=\mu^{n_0+1}V,
\]
write $J=\deg_{\mathrm{Fl}}(\mu^{n_0}V)>0$ and
$v=\sum_Zv_Z>0$, with the original component volumes $v_Z$.
The original normalized test intersection therefore satisfies
\[
 \beta_B:=\frac{\sum_Z I_{B,Z}(F_0)}{vfq(n_0+1)J}
                       \ge-\frac{107}{2}\Lambda_B.
 \tag{HA16}
\]
All initial-prime and added-prime corrections remain in these
intersections.
\end{theorem}

\begin{proof}
Fix one datum and component before taking any limit.  Write
$a_B=-52fq\Lambda_B$ for the proved all-point lower bound.
Let $A$ be the original auxiliary ample bundle, with positive
smooth metric and point-height lower bound $-K_A$.  For a fixed
positive rational number $z=a/b$, let $F_t(z)$ and $G_t(z)$ be
the original arithmetic intersections of the $(n_0+1)$st powers
of
\[
       (\mathcal L_B,\varphi_t)+zA,
                   \qquad (\mathcal L_B,\psi_t)+zA,
\]
respectively.  Equivalently use their genuine integral tensor
powers $\mathcal L_B^b\otimes A^a$ and divide the arithmetic
intersection by $b^{n_0+1}$, and the geometric degree by $b^{n_0}$.

Proposition~\ref{prop:actual-finite-fibre-nefness} gives relative
nefness at every finite place.  Proposition~\ref{prop:actual-cusp-smoothing}
makes the curvature of the second line positive when $t$ is
sufficiently small for this fixed $z$: its negative error is
absorbed by $z\omega_A$.  It is then ample on the generic fibre
and has its original conjugation-invariant smooth metric.
The point-height theorem gives, for every algebraic point and
the same sufficiently small $t$,
\[
                       h_{B,z,t}(P)\ge a_B-zK_A.
 \tag{HA17}
\]
The smoothing threshold is chosen before $P$.  Zhang's lower
intersection inequality, with precisely these finite and infinite
positivity hypotheses, gives
\[
         G_t(z)\ge(n_0+1)(a_B-zK_A)d(z),
 \qquad
 d(z)=\int_X(T+z\omega_A)^{n_0},\quad T=\theta+dd^c\varphi.
 \tag{HA18}
\]
The quoted inequality and its normalization are established in
the point-to-intersection argument\cite{height-assembly-height-record}.

The two metrics in this comparison have the same algebraic line,
finite metrics, rational factors and vertical factors.  The exact
metric-variation identity and the two energy convergence results
therefore give actual limits $F(z)$ and $G(z)$ and
\[
 G(z)-F(z)=(n_0+1)\{E_z(\psi)-E_z(\varphi)\}
       \le\frac{n_0+1}{2}\int_X U(T+z\omega_A)^{n_0}.
 \tag{HA19}
\]
The factor $1/2$ occurs only because $\psi-\varphi=U/2$;
the metric-variation formula has no additional half.
All the mixed integrals on the right are finite.  Indeed
$U=2(\psi-\varphi)$, $\psi$ is bounded, and the finite-energy
boundary theorem supplies every mixed integral weighted by
$|\varphi|$.  Polynomial interpolation in a positive auxiliary
curvature shows that $F(z)$ extends to $z=0$ with
$F(0)=A_{B,Z}$.  This is the original signed limit, not a new
positive part of that limit.

Taking $t\to0$ in (HA18), then subtracting the upper bound (HA19),
and finally taking $z\downarrow0$, gives
\[
 A_{B,Z}\ge(n_0+1)
       \left(a_B-\frac{\overline U_Z}{2}\right)d_{B,Z},
 \qquad
 \overline U_Z=\frac{\int_XUT^{n_0}}{\int_XT^{n_0}}.
 \tag{HA20}
\]
The geometric denominator is strictly positive: the original
interior curvature is the positive invariant volume density,
and the arithmetic component has positive finite volume.
The exact curvature-to-Haar dictionary, including the original
sheet multiplicities, gives
$\overline U_Z\le3fq\Lambda_B$ by the proved isotropic counting
and cusp-mean estimate\cite{height-assembly-cusp-record}.
This proves (HA15).  The order of inequalities is essential:
the bounded-metric intersection has a lower bound, and its
increase over the original metric has an upper bound.

We finish by checking that this is the actual flag test.
On the original flag model, $\mu$ is the first arithmetic Chern
class of the pullback of $\mathcal L_B$.  The determinant of a
metric flag is the product of the metrics on its successive
quotients, so this is an isometry as well as an algebraic line
identity.  It includes the original corrections at all primes.
The class $V$ has the fibre dimension $d-n_0$.
On a common projective model comparing the original flag and base
presentations, let $\delta>0$ be the original projection degree,
and let $c_{\mathrm{fib}}$ be the degree of $V$ on the generic
fibre.  The arithmetic projection formula for the smooth metrics,
and the geometric projection formula, give respectively
\[
 \delta I_{B,Z}(F_0)=c_{\mathrm{fib}}A_{B,Z},\qquad
 \delta v_ZJ=c_{\mathrm{fib}}d_{B,Z}.
 \tag{HA21}
\]
For the first equality take the just-proved metric limit after
the smooth projection formula.  Pullback of each already locally
free corrected bundle, and of its actual determinant presentation,
is exact on the common model.  Thus the vertical factors in the
two formulas are identical.  The second equality makes
$c_{\mathrm{fib}}>0$ and retains every finite projection and
component factor.  This is also the precise common-cover
calculation in the original test identification
\cite{height-assembly-source}.

Canceling $\delta$ and $c_{\mathrm{fib}}$, rather than setting
either to one, yields
\[
 I_{B,Z}(F_0)=
       \frac{A_{B,Z}}{(n_0+1)d_{B,Z}}\,(n_0+1)v_ZJ.
\]
Summing gives
$\beta_B=(fq)^{-1}\sum_Z(v_Z/v)
A_{B,Z}/((n_0+1)d_{B,Z})$, proving (HA16).
The initial corrections are already in the original determinant
presentation, while the new primes occur in its displayed
vertical product.  They have not been discarded because their
complex curvature vanishes.
\end{proof}
